\section{Matching augmentation constraints for a fixed number of ports}
\label{sec:matching-augmentation}

The protected small-cut theorem applies after adding any fixed matching between ports.
It supplies uniform inequalities for the higher terminal-state sectors.

Let $Y$ be a normalized $k$-pole of rank $h$.
For a nonempty matching $J$ between distinct ports, let $R_J(Y)$ count the following forests.
Their odd-degree vertices are exactly the endpoints of $J$.
Their path components define a pairing $\beta$ on those endpoints.
The union of $J$ and $\beta$, treated as two matchings with distinguished edge colors, must be one alternating cycle.
For $|J|=1$, the alternating cycle has two parallel colored edges, and $R_J$ is the corresponding terminal-path count.

Write $\rho_J=2^{-h}R_J(Y)$.
If $J$ has two edges, $R_J$ is the sum of the two pairing counts other than $J$.
If $J$ has three edges, $R_J$ counts pairing forests whose pairing shares no edge with $J$.

\begin{proposition}[Exact augmentation identity]
\label{prop:matching-augmentation-identity}
Let $A$ be a matching on the ports of $Y$, and choose an edge $e=ab\in A$.
Add all edges of $A\setminus\{e\}$ to $Y$ and call the resulting pole $Y_A$.
The ports $a,b$ remain unsaturated.
Then
\[
 P_{ab}(Y_A)=\sum_{\substack{J\subseteq A\\e\in J}}R_J(Y).
\]
The rank of $Y_A$ is $h+|A|-1$.
\end{proposition}

\begin{proof}
Separate terminal paths according to the subset $I\subseteq A\setminus\{e\}$ of added edges that they use.
Deleting those added edges leaves a forest of vertex-disjoint paths.
Its odd-degree vertices are exactly the endpoints of $J=I\cup\{e\}$.
Adding back $I$ joins these components into one $a$--$b$ path precisely when $J$ and its forest pairing form one alternating cycle.
This gives a bijection with the forests counted by $R_J(Y)$.
The disjoint alternatives over $I$ prove the identity.
Adding each edge increases the rank by one, giving the rank formula.
\end{proof}

\begin{theorem}[Uniform matching augmentation inequalities]
\label{thm:matching-augmentation-inequalities}
Fix $k$.
Suppose $Y$ lies inside a cubic graph with no cyclic cut of size at most four.
Uniformly over every nonempty port matching $A$ and every $e\in A$,
\begin{equation}\label{eq:matching-augmentation-bound}
 \sum_{\substack{J\subseteq A\\e\in J}}\rho_J
 \le2^{|A|-2}+o(1)
 \qquad(h\to\infty).
\end{equation}
For $|A|=1$, the right side is $1/2+o(1)$.
\end{theorem}

\begin{proof}
Apply Proposition~\ref{prop:matching-augmentation-identity}.
Protect every endpoint of the $|A|-1$ added edges.
If the augmented terminal-path density exceeded one half by a fixed amount along unbounded ranks, Theorem~\ref{thm:high-density-protected-cyclic-sets} would apply.
It would produce two disjoint cyclic sets avoiding the protected vertices and all remaining ports.
Removing the added edges and restoring the original port incidences leaves their cycles and boundaries unchanged.
This would give a cyclic cut of size at most four in the original cubic graph.

Thus
\[
 P_{ab}(Y_A)\le(1/2+o(1))2^{h+|A|-1}.
\]
Divide by $2^h$ and use the exact identity.
There are only finitely many matchings for fixed $k$, so the error is uniform.
\end{proof}

For six ports, a perfect matching $A=\{ab,ij,kl\}$ gives the inequality
\[
 a_{ab}
 +\rho_{\{ab,ij\}}
 +\rho_{\{ab,kl\}}
 +\rho_{\{ab,ij,kl\}}
 \le2+o(1).
\]
The two middle terms are noncap four-terminal forest sums.
The last term is the compatible six-terminal forest mass.
These inequalities are proved inputs to the six-cut problem.
No six-cut error theorem is asserted here.
